% =====================================================================
%  Lecture 21 -- Network Slicing                          (English edition)
% =====================================================================
\documentclass[11pt,a4paper]{article}
\providecommand\ACNroot{../..}
\input{\ACNroot/templates/lecture-en}

\ACNset{number}{21}
\ACNset{title}{Network Slicing}
\ACNset{shorttitle}{Network slicing}
\ACNset{titleother}{برش‌بندی شبکه}
\ACNset{part}{Part IV -- 5G system and procedures (Lectures 14, 19--22)}
\ACNset{time}{6--8 hours (3 h reading, 3--5 h exercises)}
\ACNset{level}{Advanced}
\ACNset{prereq}{Lectures 2 (Erlang B), 16, 19 and 20; normal distribution}
\ACNset{gen}{5}

\begin{document}
\ACNcover

% ---------------------------------------------------------------------
\section{Learning objectives}
After studying this lecture you should be able to:
\begin{objectives}
  \item \textbf{Define} a network slice and \textbf{decode} the S-NSSAI and the NSSAI sets used in signalling.
  \item \textbf{Trace} slice selection at registration and slice-specific PDU session set-up.
  \item \textbf{Explain} how slices are realised in the RAN, transport and core, and the isolation options.
  \item \textbf{Describe} slice management: lifecycle, templates and the roles of operator and tenant.
  \item \textbf{Quantify} the statistical-multiplexing cost of isolation, RAN resource partitions and slice
        admission control.
\end{objectives}

% ---------------------------------------------------------------------
\section{Prerequisite knowledge}
You need the 5GC architecture and procedures (Lectures~19--20), the virtualisation concepts of Lecture~16, and the Erlang-B formula
of Lecture~2. The multiplexing model uses the normal distribution: for a Gaussian load with mean $\mu$ and standard deviation $\sigma$,
the capacity exceeded with probability $\epsilon$ is $\mu + z_{1-\epsilon}\sigma$ ($z_{0.99} = 2.326$).

\paragraph{Key terms} The vocabulary of this lecture with its Persian equivalents:
\begin{keyterms}
Network slice & \fa{برش شبکه} & Logical network providing specific capabilities and characteristics. \\
S-NSSAI & \fa{شناسهٔ برش} & Single network slice selection assistance information (SST + SD). \\
SST / SD & \fa{نوع برش / متمایزکننده} & Slice/service type and optional slice differentiator. \\
NSSAI & \fa{مجموعهٔ S-NSSAIها} & Configured, requested, allowed and rejected sets. \\
NSSF & \fa{تابع انتخاب برش} & NF that selects allowed slices and the AMF set. \\
NSSAA & \fa{احراز هویت ویژهٔ برش} & Slice-specific authentication by a third-party AAA server. \\
NSAC & \fa{کنترل پذیرش برش} & Limits on UEs or PDU sessions per slice. \\
Isolation & \fa{جداسازی} & Degree to which one slice is protected from another. \\
Tenant & \fa{مستأجر} & Customer (e.g.\ enterprise) for whom a slice is operated. \\
Slice template & \fa{الگوی برش} & Standardised description of slice attributes (e.g.\ latency, throughput). \\
\end{keyterms}

% ---------------------------------------------------------------------
\section{Introduction and motivation}
The three IMT-2020 scenarios (Lecture~14) make conflicting demands: an eMBB service wants throughput and tolerates delay, a URLLC
service wants guaranteed latency and reliability, an mMTC service wants huge numbers of cheap, rarely active connections. Building three
separate networks would be uneconomic; running all on one network with best-effort sharing would fail the most demanding customers.
\emph{Network slicing} runs several logical networks on one physical infrastructure, each with its own capabilities, policies and,
where needed, its own network functions \cite{ts23501,foukas2017slicing}.

Slicing is also a business concept: an operator can offer a slice to an enterprise, a public-safety agency or a vertical industry as a
product with guaranteed characteristics, managed through standard interfaces \cite{ts28530}. It is made possible by the technologies of
Part~III: virtualisation creates dedicated instances cheaply, SDN programs the transport, and the SBA lets NFs be selected per slice.

% ---------------------------------------------------------------------
\section{Fundamental concepts}

\subsection{Identifying slices}
A slice is identified by an \emph{S-NSSAI} (\cref{fig:snssai}): an 8-bit \emph{slice/service type} (SST) and an optional 24-bit
\emph{slice differentiator} (SD) that distinguishes several slices of the same type (e.g.\ two enterprises' URLLC slices) \cite{ts23003}.
Standardised SST values include 1 (eMBB), 2 (URLLC), 3 (massive IoT) and 4 (V2X), with further values added in later releases;
operators may also define their own. A UE handles sets of S-NSSAIs:
\begin{itemize}
  \item \textbf{configured NSSAI}: provisioned in the UE per PLMN;
  \item \textbf{requested NSSAI}: sent in the Registration Request;
  \item \textbf{allowed NSSAI}: returned by the network for the registration area (at most 8 S-NSSAIs);
  \item \textbf{rejected NSSAI}: requested slices refused, with a cause (e.g.\ not available in this area).
\end{itemize}

\begin{acnfigure}
  \resizebox{\linewidth}{!}{\acnfig{lecture-21/snssai}}
  \caption{S-NSSAI structure and the NSSAI sets exchanged at registration.}
  \label{fig:snssai}
\end{acnfigure}

\subsection{What a slice consists of}
An end-to-end slice spans the RAN, the transport network and the core (\cref{fig:slices}). Some functions are typically shared (AMF,
NSSF, NRF, UDM, the RAN nodes), others may be slice-specific (SMF, UPF, PCF, application functions). A UE can use up to eight slices
simultaneously, but all through the same AMF; each PDU session belongs to exactly one slice.

\begin{acnfigure}
  \resizebox{\linewidth}{!}{\acnfig{lecture-21/slices}}
  \caption{Three end-to-end slices sharing RAN and common control functions, with slice-specific SMFs and UPFs.}
  \label{fig:slices}
\end{acnfigure}

\subsection{Isolation}
Isolation can be achieved at several levels: \emph{logical} (separate policies and QoS on shared NFs), \emph{virtual} (dedicated NF
instances on shared infrastructure) or \emph{physical} (dedicated hardware, spectrum or sites). Stronger isolation protects a slice from
others' load, faults and security incidents, but forfeits the statistical-multiplexing gain of sharing (\cref{sec:mux}). The right level
depends on the service-level agreement: a public-safety slice may need dedicated UPFs and reserved radio resources, while a consumer video
slice can share almost everything.

% ---------------------------------------------------------------------
\section{Detailed technical explanation}

\subsection{Slice selection}
At registration the AMF checks the requested NSSAI against the subscribed S-NSSAIs from the UDM and, if necessary, asks the NSSF, which
returns the allowed NSSAI for the UE's location and possibly a different AMF set better suited to the slices (\cref{fig:sel}). The AMF may
then re-route the registration to an AMF of that set. When the UE requests a PDU session for an S-NSSAI and DNN, the AMF discovers an SMF
serving that slice through the NRF, and the SMF selects a UPF of the slice. Slice information is also given to the RAN, which may use it
for resource allocation and mobility decisions.

\begin{acnfigure}
  \resizebox{\linewidth}{!}{\acnfig{lecture-21/selection}}
  \caption{Slice selection at registration and slice-specific PDU session set-up (simplified).}
  \label{fig:sel}
\end{acnfigure}

\subsection{Slice-specific authentication and admission control}
Release~16 added \emph{network slice-specific authentication and authorisation} (NSSAA): after primary authentication, access to a slice can
additionally depend on credentials verified by an AAA server of the slice tenant (e.g.\ an enterprise), so that the tenant controls which
devices use its slice. Release~17 added \emph{network slice admission control} (NSAC): the network can limit the number of registered UEs
and of PDU sessions per slice, enforcing the capacity sold to a tenant \cite{ts23501}.

\subsection{Slicing in the RAN and transport}
The RAN is shared; slices are realised by scheduling policies: minimum, maximum and dedicated shares of physical resource blocks per slice,
slice-aware admission control and priority, and possibly separate bandwidth parts or carriers. The transport network maps slices to
virtual private networks, segment-routing policies or hard-isolated channels with guaranteed bandwidth and latency. Without these,
a core slice alone guarantees nothing end to end.

\subsection{Slice management}
3GPP management specifications describe slices as managed objects: a \emph{network slice} composed of \emph{network slice subnets} (RAN,
transport, core), with a lifecycle of preparation, commissioning, operation and decommissioning \cite{ts28530}. Customers' requirements are
expressed with slice templates -- for example the GSMA's generic network slice template lists attributes such as throughput, latency,
availability, number of devices and isolation -- which the operator's orchestration translates into NF instances, policies and RAN and
transport configuration, using NFV MANO (Lecture~16) and SDN controllers (Lecture~15).

% ---------------------------------------------------------------------
\section{Mathematical models}

\subsection{The cost of isolation}\label{sec:mux}
\paragraph{Model} Slice $i$ has a Gaussian load with mean $\mu_i$ and standard deviation $\sigma_i$, independent of the others. To keep the
probability of overload below $\epsilon$, dedicated capacity per slice and shared capacity for all slices are
\begin{equation}\label{eq:dedicated}
  C_i = \mu_i + z\,\sigma_i, \qquad
  C_{\mathrm{shared}} = \sum_i \mu_i + z\sqrt{\sum_i \sigma_i^2},
\end{equation}
with $z = z_{1-\epsilon}$. The \emph{multiplexing gain} of sharing is
\begin{equation}\label{eq:gain}
  G = \frac{\sum_i C_i}{C_{\mathrm{shared}}} \ge 1 .
\end{equation}
\paragraph{Example} Three slices with $\mu = 100$, $\sigma = 30$ (in any capacity unit), $\epsilon = 1\,\%$ ($z = 2.326$): dedicated $3\times169.8 =
509.3$, shared $300 + 2.326\sqrt3\times30 = 420.9$; $G = 1.21$ -- isolation costs 21\,\% more capacity (\cref{fig:mux}).
\paragraph{Interpretation} The gain grows with the number of slices and with the burstiness $\sigma/\mu$; it vanishes for constant loads.
Isolation is therefore expensive exactly for bursty services -- the economic argument for sharing what can be shared and isolating only what
the SLA requires. Correlated loads (e.g.\ all slices peaking in the same hour) reduce the gain.

\begin{acnfigure}
  \acnfig{lecture-21/mux-gain}
  \caption{Multiplexing gain of a shared pool versus dedicated slice capacity (\cref{eq:gain}, $\epsilon = 1\,\%$).}
  \label{fig:mux}
\end{acnfigure}

\subsection{RAN resource partitions}\label{sec:prb}
\paragraph{Model} A cell has $N$ resource blocks; slice $i$ is guaranteed a minimum share $\rho_i$ and, with a robust modulation and coding
scheme, a rate $r_i$ per resource block. Its guaranteed throughput is
\begin{equation}\label{eq:prb}
  R_i^{\min} = \lfloor\rho_i N\rfloor\, r_i ,
\end{equation}
while unused guaranteed resources may be lent to other slices unless they are \emph{dedicated}.
\paragraph{Example} $N = 273$ (100\,MHz at 30\,kHz), URLLC minimum 20\,\% ($54$ resource blocks), $r = 0.5$\,Mbit/s per block: $R^{\min} = 27$\,Mbit/s,
sufficient for many control flows, with the rest of the cell available to eMBB when URLLC is idle.

\subsection{Slice admission control}\label{sec:nsac}
\paragraph{Model} If a tenant's slice is limited to $c$ simultaneous PDU sessions and sessions arrive as a Poisson stream with offered load
$A$ Erlang, the probability that a new session is rejected is the Erlang-B formula of Lecture~2:
\begin{equation}\label{eq:block}
  B(A, c) = \frac{A^{c}/c!}{\sum_{k=0}^{c} A^{k}/k!} .
\end{equation}
\paragraph{Example} $A = 40$\,E: with $c = 50$, $B = 1.9\,\%$; with $c = 60$, $B = 0.07\,\%$. A quota of $c = 53$ is the smallest giving $B \le 1\,\%$.
The quota is part of the contract: it protects other slices and defines what the tenant pays for.

% ---------------------------------------------------------------------
\section{Network architecture and diagrams}
\Cref{fig:slices} shows a common deployment pattern: shared RAN with per-slice scheduling policies, transport VPNs per slice, shared AMF and
subscriber data, and per-slice SMF and UPF -- the URLLC slice's UPF placed at the edge, the eMBB slice's UPFs centrally. An enterprise slice may
additionally have a UPF on the enterprise premises (Lecture~19) and NSSAA towards the enterprise's AAA server. The management system keeps the
mapping from each slice to its subnets and resources.

% ---------------------------------------------------------------------
\section{Step-by-step procedures}

\begin{procedurebox}{Registration with slice selection}\label{proc:sel}
\begin{steps}
  \item The UE sends a Registration Request with the requested NSSAI (from its configured NSSAI).
  \item The AMF retrieves the subscribed S-NSSAIs from the UDM.
  \item If it cannot decide alone, the AMF queries the NSSF with requested and subscribed NSSAI and the UE's TAI; the NSSF returns the allowed
        NSSAI and possibly a target AMF set.
  \item If another AMF set is needed, the registration is re-routed; slices requiring NSSAA are authorised with the tenant's AAA server.
  \item Registration Accept returns allowed and rejected NSSAI; the UE uses only allowed S-NSSAIs in PDU session requests.
\end{steps}
\end{procedurebox}

\begin{procedurebox}{Commissioning an enterprise slice}\label{proc:commission}
\begin{steps}
  \item The tenant's requirements are captured in a slice template (throughput, latency, devices, area, isolation, availability).
  \item The operator's management system decomposes the slice into RAN, transport and core subnets and checks feasibility.
  \item NFV orchestration instantiates slice-specific SMF/UPF (e.g.\ an on-premises UPF); SDN controllers set up transport paths; RAN policies
        reserve resource shares in the covered cells.
  \item Subscriptions of the tenant's devices are updated with the new S-NSSAI; the NSSF configuration and NSAC quotas are updated.
  \item The slice is activated and its performance is monitored against the SLA.
\end{steps}
\end{procedurebox}

% ---------------------------------------------------------------------
\section{Worked examples}

\begin{examplebox}{Decoding S-NSSAIs}\label{ex:snssai}
\textbf{Problem.} A UE's allowed NSSAI contains (SST 1), (SST 2, SD \texttt{0x00A1B2}) and (SST 2, SD \texttt{0x00C3D4}). Interpret it.

\textbf{Solution.} One default eMBB slice without differentiator and two URLLC slices distinguished by their SDs -- for example two different
enterprise or industrial customers. The UE can open PDU sessions in all three slices simultaneously through the same AMF.
\end{examplebox}

\begin{examplebox}{Isolation cost}\label{ex:mux}
\textbf{Problem.} Using \cref{eq:dedicated}, how much capacity does an operator save by sharing UPF capacity among the three slices of \cref{sec:mux}
instead of dedicating it, and what does the URLLC tenant risk?

\textbf{Solution.} $509.3 - 420.9 = 88.4$ units (17\,\% of the dedicated total). The URLLC slice then competes with the others during joint
peaks; if its SLA requires guaranteed performance independent of others' traffic, its share must be reserved (logical isolation with
priority) or dedicated (virtual isolation).
\end{examplebox}

\begin{examplebox}{Dimensioning a quota}\label{ex:nsac}
\textbf{Problem.} A logistics company expects 40\,E of simultaneous sessions in its slice and accepts 1\,\% blocking. What NSAC quota should be
configured and how does the blocking change if its load grows by 12.5\,\%?

\textbf{Solution.} $c = 53$ (\cref{sec:nsac}). At $A = 45$\,E with $c = 50$ the blocking would be 5.4\,\%; with $c = 53$ it is 3.1\,\% -- still well above
1\,\% -- the quota (and price) must be renegotiated as the tenant grows.
\end{examplebox}

% ---------------------------------------------------------------------
\section{Technology comparison}
\begin{table}[htbp]
\centering\small
\caption{Ways of differentiating service in mobile networks.}\label{tab:cmp}
\rowcolors{2}{acnmist}{white}
\begin{tabularx}{\linewidth}{@{}l Y Y Y@{}}
\toprule
\rowcolor{white} & QoS classes (4G/5G) & Dedicated core (DECOR, 4G) & Network slicing (5G) \\
\midrule
Granularity    & Per flow/bearer                & Per UE (core selection)      & Per PDU session, end to end \\
Scope          & Packet treatment               & Core NFs                     & RAN, transport, core, management \\
Isolation      & None (priorities)              & Core instances               & Logical to physical \\
Per-UE use     & Many flows                     & One dedicated core           & Up to 8 slices \\
Customer view  & Not exposed                    & Not exposed                  & Product with SLA, tenant management \\
\bottomrule
\end{tabularx}
\end{table}

\paragraph{Slicing and QoS} Slices and QoS flows (Lecture~22) are complementary: a slice defines \emph{which} network and resources serve a
session; QoS flows define how packets \emph{within} it are treated. A URLLC slice still uses QoS flows with appropriate 5QIs.

% ---------------------------------------------------------------------
\section{Key takeaways}
\begin{takeawaybox}
\begin{itemize}
  \item A slice is an end-to-end logical network identified by S-NSSAI = SST (8 bits) + optional SD (24 bits).
  \item NSSAI sets: configured, requested, allowed (max.\ 8), rejected; NSSF selects slices and AMF sets.
  \item Slices combine shared and dedicated NFs; RAN and transport must also be sliced for end-to-end guarantees.
  \item Isolation costs multiplexing gain: $G = \sum(\mu_i + z\sigma_i)/(\sum\mu_i + z\sqrt{\sum\sigma_i^2})$.
  \item RAN partitions give $R_i^{\min} = \lfloor\rho_i N\rfloor r_i$; NSAC quotas are dimensioned with Erlang B.
  \item NSSAA and NSAC let tenants control access and operators enforce contracted capacity.
\end{itemize}
\end{takeawaybox}

% ---------------------------------------------------------------------
\section{Exercises}

\begin{exercisebox}{1}{Slices or QoS?}\label{exr:slqos}
For each requirement, say whether a separate slice or a QoS flow within an existing slice is the better tool: (a) prioritising video calls
over downloads for consumers; (b) a police network that must keep working when the public network is congested; (c) a factory whose data must
not leave its premises.
\end{exercisebox}

\begin{exercisebox}{1}{Isolation levels}\label{exr:iso}
Give an example of logical, virtual and physical isolation in the core and in the RAN.
\end{exercisebox}

\begin{exercisebox}{2}{Bursty slices}\label{exr:gain}
Repeat the calculation of \cref{sec:mux} with $\sigma = 60$. What is the gain, and what does it mean for slices carrying bursty traffic?
\end{exercisebox}

\begin{exercisebox}{2}{RAN share}\label{exr:prb}
A URLLC slice is guaranteed 15\,\% of 273 resource blocks with a robust MCS of 0.2\,Mbit/s per block. Compute its guaranteed rate. How many 64\,kbit/s
control streams fit, and what happens to unused blocks if the share is ``minimum'' rather than ``dedicated''?
\end{exercisebox}

\begin{exercisebox}{3}{Quota design}\label{exr:nsac}
Show that the smallest quota with $B \le 1\,\%$ at $A = 40$\,E is $c = 53$ (use the recursion $B(A,k) = A B(A,k-1)/(k + A B(A,k-1))$). Then estimate the quota
needed for $A = 80$\,E and compare the ratio $c/A$; explain the trend.
\end{exercisebox}

\begin{exercisebox}{3}{Designing a slice}\label{exr:design}
A hospital wants a slice for 200 connected devices (monitors, imaging) with 10\,ms latency, 99.99\,\% availability inside its campus, and data that must
stay on site. Propose the S-NSSAI, the shared and dedicated components, the UPF placement, RAN measures and admission/authorisation features.
\end{exercisebox}

\subsection*{Solutions}
\addcontentsline{toc}{subsection}{Solutions}

\begin{solutionbox}{\ref{exr:slqos}}
(a) QoS flows (different 5QIs) in the consumer slice. (b) A slice with reserved RAN resources, priority and preemption (and possibly dedicated core NFs)
-- isolation from public load is the requirement. (c) A slice with an on-premises UPF and local breakout (possibly NSSAA), since the requirement is about
where the user plane terminates, not packet priority.
\end{solutionbox}

\begin{solutionbox}{\ref{exr:iso}}
Core: logical -- shared SMF/UPF with per-slice policies; virtual -- dedicated SMF/UPF instances on shared servers; physical -- dedicated servers or site.
RAN: logical -- per-slice priority in a shared scheduler; virtual -- reserved/dedicated PRB shares or a separate bandwidth part; physical -- a dedicated
carrier or dedicated small cells.
\end{solutionbox}

\begin{solutionbox}{\ref{exr:gain}}
Dedicated $3\times(100 + 2.326\times60) = 718.7$; shared $300 + 2.326\sqrt3\times60 = 541.7$; $G = 1.33$. The burstier the traffic, the more capacity
isolation wastes: bursty best-effort slices should share; isolation should be reserved for slices whose SLA demands it.
\end{solutionbox}

\begin{solutionbox}{\ref{exr:prb}}
$\lfloor0.15\times273\rfloor = 40$ blocks (41 if rounded), $R^{\min} = 40\times0.2 = 8$\,Mbit/s: 125 streams of 64\,kbit/s (fewer with headroom). With a
\emph{minimum} share, unused blocks are lent to other slices and reclaimed when URLLC traffic appears (possible short scheduling delay); with a
\emph{dedicated} share they stay idle, giving the strongest guarantee at a capacity cost.
\end{solutionbox}

\begin{solutionbox}{\ref{exr:nsac}}
The recursion gives $B(40,52) = 1.1\,\%$ and $B(40,53) = 0.82\,\%$, hence $c = 53$ ($c/A = 1.33$). For $A = 80$\,E the recursion gives $c = 96$ ($c/A = 1.2$):
larger pools need relatively less headroom -- the same economy of scale as in trunk dimensioning (Lecture~2) and the multiplexing gain of \cref{sec:mux}.
\end{solutionbox}

\begin{solutionbox}{\ref{exr:design}}
S-NSSAI: SST 2 (URLLC) or an operator-specific SST, with an SD identifying the hospital. Shared: AMF, NSSF, UDM. Dedicated: SMF and a UPF on the hospital campus
(local breakout, data on site), possibly redundant UPFs for 99.99\,\%. RAN: indoor small cells with a minimum or dedicated PRB share and slice priority; latency
via 30\,kHz numerology or mini-slots. Access: NSSAA with the hospital's AAA server; NSAC quota of about 200 UEs; transport path with guaranteed bandwidth. Monitoring
of latency and availability against the SLA.
\end{solutionbox}

% ---------------------------------------------------------------------
\section{Further reading}
Slicing in the 5G system is specified in TS~23.501 \cite{ts23501} (clause 5.15) with procedures in TS~23.502 \cite{ts23502}; S-NSSAI encoding in TS~23.003
\cite{ts23003}; management concepts in TS~28.530 \cite{ts28530}; service requirements in TS~22.261 \cite{ts22261}. Foukas et al.\ \cite{foukas2017slicing} survey
slicing architectures and research questions; Rommer et al.\ \cite{rommer2020core} cover slicing in the 5G core.

\printbibliography[heading=subbibliography,title={\ACNstrRefs}]
\end{document}
